\chapter{Modos de Operación y Procedimientos}
\label{ch:modos-operacion}

Este capítulo describe los modos de operación del \textbf{CMOS Microscopy Payload
(CMOSM)}, la lógica de transición entre estados y los procedimientos asociados a
la adquisición de imágenes en microscopía sin lentes en configuración de
contacto. El enfoque prioriza la operación nominal del sistema y documenta los
modos experimentales implementados para evaluación computacional.

% -------------------------------------------------------------------
\section{Descripción general}

El payload CMOSM opera mediante una máquina de estados bien definida que controla
la inicialización del hardware, la visualización preliminar, la captura de
imágenes, el registro de metadatos y el apagado seguro del sistema.

El \textbf{modo nominal de operación} corresponde a la captura de una imagen
única bajo parámetros de iluminación y adquisición controlados, siendo suficiente
para la identificación visual de estructuras celulares en muestras biológicas
planas. De manera complementaria, el sistema implementa \textbf{modos
experimentales} orientados a la evaluación de técnicas de reconstrucción
computacional, como Fourier Ptychographic Microscopy (FPM).

La lógica de control es ejecutada por software en la Raspberry Pi, actuando como
computadora de carga útil y comunicándose con el sistema de OBDH cuando es
requerido.

% -------------------------------------------------------------------
\section{Modos operativos del payload}

\begin{table}[H]
\centering
\begin{tabular}{p{3.5cm} p{9.5cm}}
\toprule
\textbf{Modo} & \textbf{Descripción} \\
\midrule
Idle & Estado pasivo sin captura activa; hardware energizado. \\
Init & Inicialización del sensor CMOS y del sistema de iluminación OLED. \\
Preview & Visualización en tiempo real y ajuste de parámetros de captura. \\
Capture-Single & \textbf{Modo nominal}. Captura de una imagen única bajo
condiciones controladas. \\
Capture-FPM & \textbf{Modo experimental}. Captura secuencial con iluminación
estructurada para evaluación de FPM. \\
Metadata & Generación y almacenamiento de metadatos asociados a la captura. \\
Diagnostics & Evaluación rápida de calidad de imagen (SNR, saturación, estabilidad). \\
Safe & Apagado seguro del sistema y liberación de recursos. \\
\bottomrule
\end{tabular}
\caption{Modos operativos del CMOS Microscopy Payload.}
\end{table}

% -------------------------------------------------------------------
\section{Procedimiento de captura nominal (Capture-Single)}

El modo \textit{Capture-Single} constituye el \textbf{modo nominal de operación
del payload}. Permite la adquisición de una imagen única bajo parámetros de
captura fijos y controlados, siendo suficiente para la caracterización visual de
estructuras celulares en muestras biológicas planas.

\subsection*{Secuencia de operación}

\begin{enumerate}
    \item Inicialización del sensor CMOS y del sistema de iluminación OLED.
    \item Configuración manual de exposición y ganancia del sensor.
    \item Aplicación del patrón de iluminación seleccionado.
    \item Verificación del histograma para evitar saturación.
    \item Captura de imagen en formato sin compresión.
    \item Almacenamiento de la imagen.
    \item Generación y almacenamiento de metadatos asociados.
\end{enumerate}

Este modo se utiliza para operación científica nominal, calibración y
caracterización de muestras.

% -------------------------------------------------------------------
\section{Parámetros de adquisición y metadatos}

Durante cada captura, el sistema registra automáticamente un conjunto de
parámetros de adquisición y métricas básicas de calidad, los cuales se almacenan
junto con la imagen adquirida para asegurar trazabilidad experimental.

\begin{table}[H]
\centering
\begin{tabular}{p{4cm} p{3cm} p{7cm}}
\toprule
\textbf{Parámetro} & \textbf{Unidad} & \textbf{Descripción} \\
\midrule
ExposureTime & $\mu$s & Tiempo de integración del sensor CMOS. \\
AnalogueGain & -- & Ganancia analógica aplicada al sensor. \\
FrameDuration & $\mu$s & Duración del frame del sensor. \\
Resolution & píxeles & Resolución de captura (2592$\times$1944). \\
OLED Pattern ID & -- & Identificador del patrón de iluminación utilizado. \\
OLED Intensity & -- & Nivel relativo de intensidad del OLED. \\
Timestamp & -- & Marca temporal de la captura. \\
Histogram Mean & ADU & Valor medio del histograma de la imagen. \\
Saturation Ratio & \% & Porcentaje de píxeles saturados. \\
Sample Tag & -- & Identificador experimental de la muestra. \\
\bottomrule
\end{tabular}
\caption{Parámetros registrados por captura en el CMOSM.}
\end{table}

% -------------------------------------------------------------------
\section{Modo experimental Capture-FPM}

El modo \textit{Capture-FPM} permite la adquisición secuencial de imágenes bajo
iluminación estructurada con el objetivo de evaluar técnicas experimentales de
reconstrucción computacional, como Fourier Ptychographic Microscopy (FPM). Este
modo no es requerido para la operación nominal del payload.

\subsection*{Secuencia experimental}

\begin{enumerate}
    \item Fijar parámetros de exposición y ganancia.
    \item Ejecutar un barrido de patrones OLED predefinidos.
    \item Capturar una imagen por patrón.
    \item Registrar metadatos individuales y globales de la secuencia.
\end{enumerate}

El conjunto de imágenes obtenido se utiliza exclusivamente para análisis
experimental y evaluación de algoritmos de reconstrucción.

% -------------------------------------------------------------------
\section{Referencia de implementación en software}

A modo de referencia, el Listado~\ref{lst:capture-single} muestra un fragmento
simplificado del flujo lógico empleado para la captura nominal. Este listado se
incluye únicamente con fines descriptivos y no constituye una especificación de
implementación.

\begin{algorithm}[H]
\caption{Ejemplo de flujo de captura nominal}
\label{lst:capture-single}
\begin{algorithmic}[1]
\State Desactivar auto-exposición y auto-balance de blancos
\State Configurar \texttt{ExposureTime} y \texttt{AnalogueGain}
\State Aplicar patrón OLED seleccionado
\State Capturar imagen desde el sensor CMOS
\State Calcular métricas básicas de histograma
\State Guardar imagen y metadatos
\end{algorithmic}
\end{algorithm}

% -------------------------------------------------------------------
\section{Modo Safe}

El modo \textit{Safe} asegura la integridad del sistema ante fallas críticas o
comandos de apagado, apagando el sistema de iluminación OLED, deteniendo la
cámara CMOS, almacenando logs de operación y liberando recursos de hardware y
software.
